% SOSP '26 TensorDex Artifacts — reviewer walkthrough
% Build:  make appendix   (two pdflatex passes)
\documentclass[11pt]{article}

\usepackage[letterpaper, margin=1in]{geometry}
\usepackage[T1]{fontenc}
\usepackage[utf8]{inputenc}
\usepackage{lmodern}
\usepackage[scaled=0.92]{helvet}
\renewcommand{\familydefault}{\sfdefault}
\usepackage{booktabs}
\usepackage{colortbl}
\usepackage{graphicx}
\usepackage{enumitem}
\usepackage{xcolor}
\usepackage{titlesec}
\usepackage{parskip}
\usepackage[colorlinks=true, linkcolor=linkblue, urlcolor=linkblue,
            pdftitle={SOSP'26 TensorDex Artifacts}]{hyperref}
\usepackage{tcolorbox}
\tcbuselibrary{listings, skins, breakable}

% ── palette ─────────────────────────────────────────────────────────
\definecolor{ink}{HTML}{202124}        % body text
\definecolor{muted}{HTML}{5F6368}      % secondary text / box labels
\definecolor{linkblue}{HTML}{1155CC}
\definecolor{shellbg}{HTML}{F1F3F4}    % shell blocks
\definecolor{termbg}{HTML}{FBF3DB}     % terminal-output blocks
\definecolor{cmdred}{HTML}{C0392B}     % comments / keywords, like the ref doc
\definecolor{strgreen}{HTML}{1E8E3E}
\definecolor{noticeblue}{HTML}{1A0DAB}
\color{ink}

% ── headings: big, flush-left, sans (Google-Docs-guide look) ────────
\titleformat{\section}{\LARGE\bfseries\color{ink}}{}{0em}{}
\titlespacing*{\section}{0pt}{1.6em}{0.4em}
\titleformat{\subsection}{\Large\bfseries\color{ink}}{}{0em}{}
\titlespacing*{\subsection}{0pt}{1.2em}{0.3em}
\titleformat{\subsubsection}{\large\bfseries\color{muted}}{}{0em}{}
\titlespacing*{\subsubsection}{0pt}{1em}{0.2em}
\setlist{nosep, leftmargin=1.5em}
\pagestyle{plain}

% ── code boxes ──────────────────────────────────────────────────────
\lstdefinestyle{shellstyle}{
  language=bash,
  basicstyle=\ttfamily\small\color{ink},
  keywordstyle=\color{cmdred},
  commentstyle=\color{cmdred},
  stringstyle=\color{strgreen},
  showstringspaces=false,
  breaklines=true,
  breakatwhitespace=true,
  columns=fullflexible,
  keepspaces=true,
  upquote=true,
}
\lstdefinestyle{termstyle}{
  basicstyle=\ttfamily\scriptsize\color{ink!85},
  breaklines=true,
  columns=fullflexible,
  keepspaces=true,
  upquote=true,
}
\newtcblisting{shell}{
  listing only, listing style=shellstyle,
  enhanced, breakable,
  colback=shellbg, colframe=shellbg,
  arc=3mm, boxrule=0pt,
  left=4mm, right=3mm, top=1mm, bottom=2mm,
  title={Shell}, fonttitle=\scriptsize\sffamily,
  coltitle=muted, colbacktitle=shellbg,
  attach title to upper={\par},
  before title={}, toptitle=2mm, bottomtitle=0mm,
}
\newtcblisting{term}{
  listing only, listing style=termstyle,
  enhanced, breakable,
  colback=termbg, colframe=termbg,
  arc=2mm, boxrule=0pt,
  left=3mm, right=3mm, top=1mm, bottom=2mm,
  title={Output}, fonttitle=\scriptsize\sffamily,
  coltitle=muted, colbacktitle=termbg,
  attach title to upper={\par},
  toptitle=2mm, bottomtitle=0mm,
}
\newcommand{\code}[1]{\texttt{#1}}
\newcommand{\notice}[1]{{\color{noticeblue}\bfseries NOTICE: #1\par}}

\begin{document}

% ── title page ──────────────────────────────────────────────────────
\begin{center}
  \includegraphics[width=0.46\linewidth]{../../assets/tensordex-logo.png}\par
  \vspace{0.7em}
  {\LARGE\bfseries SOSP'26 Artifact Evaluation\par}
  \vspace{0.4em}
  {\color{muted}for \emph{TensorDex: A Compact, Tensor-Centric Storage
   System for Modern AI Models} (Paper \#323)\par}
\end{center}
\vspace{0.3em}

\begin{tcolorbox}[enhanced, colback=shellbg, colframe=shellbg, arc=3mm,
  boxrule=0pt, left=4mm, right=4mm, top=3mm, bottom=3mm]
\renewcommand{\arraystretch}{1.35}
\begin{tabular}{@{}p{0.13\linewidth}p{0.82\linewidth}@{}}
\textbf{Code} & \url{https://github.com/Antlera/TensorDex} ---
  step-by-step instructions in \code{ae/README.md}; reproduced figures
  collect in \code{ae/RESULTS.md} \\
\textbf{Data} & Hugging Face dataset \code{tensordex/tensordex-ae-cache}
  ($\sim$8.3\,GB, public), fetched by \code{make ae-cache}. The paper's
  full trace covers $\sim$40\,TB of model data and takes days to run, so we
  ship its pre-computed results; they are verifiable by re-deriving random
  samples from the raw tensor bytes we include (Tier 1) \\
\textbf{Machine} & any x86\_64 Linux box: $\geq$8\,GB RAM,
  $\geq$15\,GB disk, no GPU. Please comment on HotCRP with your SSH public
  key and a time window if you would like us to provision the paper's eval
  box (AWS c6a.48xlarge) for you \\
\textbf{Time} & kick-the-tires $\sim$5 min $\cdot$ full offline
  reproduction $\sim$45 min (\code{make reproduce-all}) \\
\end{tabular}
\end{tcolorbox}

\section{Artifact at a Glance}

Each claim in the paper maps to a single command that prints the measured
value, the expected value, and an explicit PASS/FAIL verdict. The rest of
this document walks through them: every command can be pasted into a
terminal, and every output box shows the real output from our test machine.

\begin{center}
\small
\renewcommand{\arraystretch}{1.25}
\begin{tabular}{@{}p{0.40\linewidth}p{0.24\linewidth}p{0.29\linewidth}@{}}
\toprule
\rowcolor{shellbg}
\textbf{Claim (paper)} & \textbf{Command} & \textbf{Expected} \\
\midrule
All figures derive from the published cache (Figs.~1, 2, 4, 11--16, Tab.~3) &
\code{make figures} & 38/38 charts rendered \\
\addlinespace[2pt]
70.5\,\% (FM++) / 65.1\,\% (TensorX) reduction is genuine (Figs.~1, 11) &
\code{make ae-fmpp \&\&} \code{make verify} & sampled rows re-derived from
raw bytes, bit-exact \\
\addlinespace[2pt]
Sketch-based base selection matches exact search (Fig.~12a) &
\code{make verify-recall} & Recall@1 $= 1.000$ \\
\addlinespace[2pt]
Reduction prediction error $\sim$1\,\% (Fig.~13) &
\code{make verify-predict} & held-out MAE 1.11\,\%, Pearson 99.3\,\% \\
\addlinespace[2pt]
Pipeline is functional \& lossless (\S5) &
\code{make full} & byte-exact round trip \\
\addlinespace[2pt]
FlexSplit $\approx$ ILP optimum, near-constant time (Fig.~14) &
\code{make bench-fig14} & solver sweep re-run live \\
\addlinespace[2pt]
Codec throughput (Table~3, Fig.~1-right) &
\code{make bench-table3-real}; \code{make bench-fmpp-real} & TensorX and
FM++ parallel encode/decode, byte-exact \\
\bottomrule
\multicolumn{3}{@{}l@{}}{\footnotesize Throughput variants:
\code{bench-table3} $\cdot$ \code{bench-fmpp} $\cdot$
\code{bench-baselines} (ZipNN, OpenZL,
official BitX).} \\
\end{tabular}
\end{center}

\section{Installing TensorDex}

TensorDex is a Python package with a Rust extension; it needs Rust
$\geq 1.78$, Python $\geq 3.8$, and GNU make. From a fresh clone:

\begin{shell}
git clone --branch sosp26-ae https://github.com/Antlera/TensorDex.git
cd TensorDex
# Rust toolchain, if you don't already have one:
curl --proto '=https' --tlsv1.2 -sSf https://sh.rustup.rs | sh -s -- -y
source "$HOME/.cargo/env"
# a virtualenv (stock Ubuntu ships neither venv nor a usable system pip):
sudo apt-get install -y python3-venv   # once; or use uv instead
python3 -m venv .venv && source .venv/bin/activate
# optional, saves ~5 GB: CPU-only torch (no GPU is ever used)
pip install torch --index-url https://download.pytorch.org/whl/cpu
pip install .                          # builds the Rust extension
pip install -r ae/requirements-ae.txt  # figure + verification deps
\end{shell}

The repo also ships a Dockerfile with everything (including the optional
FM++ codec) preinstalled:

\begin{shell}
docker build -t tensordex-ae .
docker run --rm tensordex-ae     # runs `make check` inside the container
docker run -it --rm -v tensordex-cache:/tensordex/ae/cache tensordex-ae bash
\end{shell}

\section{Kick the Tires ($\sim$5 minutes, no download)}

\code{make check} confirms the extension built correctly: it cross-checks the
XXH3-128 content hash against an independent library (the ids in the shipped
cache are these hashes) and runs the offline unit tests.

\begin{shell}
make check
\end{shell}
\begin{term}
OK  build + XXH3-128 hash
..................                                            [100%]
18 passed, 1 skipped in 2.78s
OK  build verified -- no download needed for this check
\end{term}

\code{make full} then runs the whole pipeline on a synthetic base + fine-tune
pair --- ingest, tensor-level dedup, delta compression, and an
integrity-verified pull --- entirely offline:

\begin{shell}
make full
\end{shell}
\begin{term}
---- 1. ingest two models (tensor-level dedup) -------------------------
  ingested base      (3 params)
  ingested finetune  (3 params)
  -> 2 models, 6 logical params, but only 4 unique tensors stored
     (embed + norm are shared -> stored once)

---- 2. compress the fine-tune against the base (raw -> delta) ---------
  attn.weight: 1.00 MiB -> 0.10 MiB (10.0x) via tensorx delta

---- 3. the delta-base graph lives in SQL (tensor_deltas) --------------
  cd2a67cdf9de...  --delta(tensorx)-->  b35165b4faf5...
  protected bases (gc-safe): ['b35165b4faf5...']

---- 4. pull the fine-tune back, with integrity verification -----------
  pull --verify -> 3 tensors, byte-exact round-trip: True
  (attn.weight was delta-decoded against the base, then re-hashed)

---- 5. sharded pull (bounds peak memory; HF-style index) --------------
  --max-shard-size 8MiB -> 3 shards + index.json

---- 6. gc keeps the base alive because a delta still needs it ---------
  removed 'base' model; gc deleted 0 tensors, protected 1 delta base(s)
  finetune still reconstructs after gc: True

demo complete [ok]
\end{term}

If both commands pass, the build is good --- everything that follows is just
data + compute.

\section{Getting the Data}

The evaluation cache lives in a public Hugging Face dataset (no token
needed): \code{results.db} --- the 11.4\,M-pair compression cache every
figure is drawn from --- plus the raw tensor blobs for Tier-1 verification
and the staged chart inputs.

\begin{shell}
make ae-cache            # everything (~8.3 GB) -> ae/cache/
make ae-cache-figures    # figures-only alternative (~6 GB, skips the blobs)
\end{shell}

This populates:

\begin{term}
ae/cache/
  results.db                       11.4M-pair compression cache (5.6 GB)
  sample_blobs/xx/yy/<id>.safetensors   raw tensors for Tier-1 verify (~2 GB)
  data/tensordb_s3/metadata.db     tensor sizes + model->tensor map
  model_hub_crawl/  tests/output/  compression_data/  ...   chart inputs
\end{term}

\section{Reproducing Experiment Results}

\notice{Throughput/QPS panels (Fig.~1 right, Fig.~12b/c, Table~3) replot
values measured on the paper's evaluation machine (c6a.48xlarge); absolute
throughput on your box will differ --- re-measure it with
\code{make bench-table3-real} (below). Every ratio-bearing result below is
re-derived, not replotted.}

\subsection{Every figure from the cache (Tier 0)}

\begin{shell}
make figures        # ~10 min -> ae/figures/*.png   (FMT=pdf for paper-quality)
\end{shell}
\begin{term}
Rendering 38 charts from ae/cache/results.db

  ok   algo_bench_q_proj.png
  ok   algo_bench_v_proj.png
  ok   bcs_memory.png
  ...
  ok   zipllm_real_bar.png
  ok   zipllm_real_cdf.png

============================================================
rendered 38 figures -> ae/figures
\end{term}

\code{ae/RESULTS.md} lays the figures out against the paper's evaluation
section, and \code{ae/FIGURE\_MAP.md} maps every chart id to its paper
figure. \code{make report} builds a self-contained single-file
\code{ae/results.html} with every figure embedded (\code{make serve} hosts
it on \code{localhost:8000}).

Two spot-checks against the paper --- Fig.~1 (left): TensorDex stores the
corpus at $0.29\times$ (70.5\,\% saved), below every baseline; and Fig.~13:
predicted vs.\ real reduction ratio hugs the diagonal:

\begin{center}
\includegraphics[height=4.1cm]{../figures/codec_storage_reduction.png}%
\hspace{0.02\linewidth}%
\includegraphics[height=4.1cm]{../figures/pred_vs_real_ratio.png}
\end{center}

\subsection{Bit-exact sample verification (Figs.~1 and 11, Tier 1)}

The cache has 11.4\,M rows; rather than recompress all of them, you draw a
random sample --- \emph{with a seed you choose} --- from the $\sim$342 rows
the shipped raw-blob bundle covers, and every sampled row is
re-derived from the raw tensor bytes shipped in the cache: the content ids
are re-hashed (XXH3-128) and the TensorX / FM++ compression ratios are
recomputed and compared against the cached values, bit for bit.

First build the optional FM++ codec (the codec behind the 70.5\,\% result; vendored
FM-Delta lib, x86\_64-linux --- skip this line and \code{verify} still checks
TensorX):

\begin{shell}
make ae-fmpp             # maturin develop --release --features fmpp
make verify              # N=200 default; or:
python ae/verify_sample.py --n 300 --seed 7     # your own seed and size
make verify-fmpp-roundtrip  # synthetic FM++ round trip; no cache needed
\end{shell}
\begin{term}
Scanning blob store ae/cache/sample_blobs ...
  326 tensor blobs available
  342 cached pairs fully covered by the blob sample

Verifying 200 random pairs (seed=0, codecs: TensorX (level 1) + FM++)
  [  1/200] OK  9a91c437e6a6...  id=[ok]  tratio got=0.345300830 exp=0.345300830
  [  2/200] OK  09fef6d3451d...  id=[ok]  tratio got=0.336914062 exp=0.336914062
  [  5/200] OK  72034a8240d1...  id=[ok]  tratio got=0.421907902 exp=0.421907902
                                          fratio got=0.359875679 exp=0.359875679
  [  6/200] OK  1d91936611eb...  id=[ok]  tratio got=0.068219866 exp=0.068219866
                                          fratio got=0.034877232 exp=0.034877232
  ... (194 more)

================================================================
content-id checks : 200/200 exact
TensorX ratio     : 199/200 bit-exact (99.5%)  (1 legacy rows)
FM++ ratio        : 94/94 bit-exact (100.0%)

RESULT: PASS  cache reproduced from raw bytes
================================================================
\end{term}

Because the cache is keyed by content hash, a matching id ties the row to
real tensor bytes, and a matching ratio proves the cached number is what the
codec actually produces on those bytes. You pick the seed, so we could not
have pre-selected favorable rows. (The one non-bit-exact row above is a
documented \emph{legacy row} --- it predates the final codec build and is
reported with its timestamp; its content ids still match, and the run still
PASSes. See Notes.) The separate FM++ round-trip command needs no cache: it
generates deterministic synthetic base/target tensors, creates fresh deltas,
decodes them, and requires byte-exact reconstruction.

\subsection{TensorSketch Recall@1 (Fig.~12a)}

Re-runs base selection twice over fingerprints recomputed from the shipped
blobs --- exact brute-force search (ground truth) vs.\ the HNSW index
TensorDex uses --- and reports how often they pick the same delta base:

\begin{shell}
make verify-recall
\end{shell}
\begin{term}
Recomputing BCS fingerprints (d=2, w=1024) for 326 tensors ...
Greedy base-selection: brute-force (ground truth) vs hnswlib (HNSW) ...

================================================================
tensors 326  .  bases 67  .  attached 259
TensorSketch Recall@1 (HNSW base == brute-force base) : 1.0000

RESULT: PASS  matches the paper's ~1.0 recall (>=0.99)
================================================================
\end{term}

\subsection{Reduction-ratio prediction (Fig.~13)}

This is reproduced \emph{as an experiment, not a replot}: TensorPred's hybrid
model is re-fit from the cached \code{(bcs\_dist, ratio)} pairs by ordinary
least squares --- no stored coefficients --- trained on a random half and
evaluated on the held-out half:

\begin{shell}
make verify-predict
\end{shell}
\begin{term}
  pairs   : 5,768,502  (bcs_dist > 0, aratio not null, bytes_in > 100KB)

  recovered hybrid coefficients (train split):
    c0=-6.0383  c1=+0.1553  c2=+0.6240  c3=+0.0135

                               MAE    Med.AE     Pearson
  held-out (50% test)        1.11%     0.79%       99.3%   (n=2,884,251)
  paper (Fig 13)             1.11%     0.79%       99.3%

  provenance vs stored pred_ratio: max|delta|=0.0001  Pearson=100.000%

  PASS: held-out prediction MAE 1.11% (< 1.5%), Pearson 99.3% (> 99%)
\end{term}

The recovered model also matches the cache's stored \code{pred\_ratio}
column to $10^{-4}$ --- confirming that column \emph{is} this model, not a
hand-tuned constant.

\subsection{FlexSplit vs.\ ILP / Primal--Dual (Fig.~14)}

Actually runs the three facility-location solvers over the cached pairwise
ratios, sweeping model counts, and re-renders both panels. Expect FlexSplit
within a few points of the ILP optimum at near-constant time, while ILP time
grows super-linearly.

\begin{shell}
make bench-fig14    # 15-25 min; ILP needs gurobipy + a (free academic)
                    # Gurobi license -- skipped automatically if absent
\end{shell}

\subsection{Codec throughput (Table 3, Fig.~1-right)}

Re-measures TensorX with the paper's exact setup: the real
Qwen2.5-7B\,/\,-Instruct pair, 2\,MB sub-chunks, zstd level~1, pure codec
compute across all cores. Below is this command's output on a c6a.48xlarge
during packaging --- reduction matches the paper exactly and throughput lands
on the published numbers. Baseline rows in Table~3 cite each system's own
published figures (ZipLLM, ZipNN, OpenZL); \code{ae/bench\_baselines.py}
re-runs the external tools. The Docker image includes the submission-time
versions, ZipNN~0.5.3 and OpenZL~v0.1.0.

\begin{shell}
make bench-table3-real   # ~30 GB model download on first run, ~32 GB RAM
make bench-table3        # synthetic variant, no download
make bench-fmpp-real     # native tensor-pair parallel FM++ + decoder
make bench-fmpp          # synthetic FM++ variant, no download
make bench-baselines     # ZipNN 0.5.3 / OpenZL v0.1.0
\end{shell}
\begin{term}
TensorX codec throughput - real model pair (14525 MB across 7481
chunks, 192 threads, zstd L1)
  reduction    0.406x (59.4% saved)   round-trip byte-exact [ok]
  compress     avg   24.3 GB/s   peak   24.7 GB/s
  decompress   avg   28.6 GB/s   peak   29.2 GB/s

paper reference (c6a.48xlarge, 192 vCPU, real Qwen2.5-7B pair):
  22.9 / 28.4 GB/s @ 59.4%
RESULT: PASS  round-trip byte-exact; throughput measured
\end{term}

The FM++ benchmark keeps every arithmetic-coded tensor stream intact and uses
Rayon to schedule independent tensor pairs across cores, matching the original
aggregate-throughput methodology. Before timing, it materializes every delta,
decodes it against its base, and requires byte-exact reconstruction. The real
variant uses the same Qwen2.5-7B Base\,/\,Instruct pair and needs roughly
40~GB of RAM. Aggregate throughput depends on the CPU count and the tensor-size
distribution; a single model pair can exhibit a long tail from its largest
embedding tensors.

\subsection{Everything in one command}

After \code{make ae-cache}, this runs all of the above (except the solver
benchmark) plus the HTML report, back to back:

\begin{shell}
make reproduce-all       # Tier 0 + Tier 1 + Tier 2 demo + report, ~30-45 min
\end{shell}

\section{Running on Your Own Models}

The same pipeline works on live Hugging Face repos --- pick a base and a
fine-tune and TensorDex reports the achieved storage reduction:

\begin{shell}
python ae/run_full.py --mode hf \
    --models Qwen/Qwen2.5-0.5B Qwen/Qwen2.5-0.5B-Instruct
python ae/render.py --only reduction_global --format pdf   # re-render one chart
python ae/run_full.py --recipe   # the full 40 TB / 2,890-model trace recipe
\end{shell}

\end{document}
